%!TEX root = deepqnet.tex

So far, we have performed experiments on seven popular ATARI games -- Beam Rider, Breakout, Enduro, Pong, Q*bert, Seaquest, Space Invaders. We use the same network architecture, learning algorithm and hyperparameters settings across all seven games, showing that our approach is robust enough to work on a variety of games without incorporating game-specific information. While we evaluated our agents on the real and unmodified games, we made one change to the reward structure of the games during training only.  Since the scale of scores varies greatly from game to game, we fixed all positive rewards to be $1$ and all negative rewards to be $-1$, leaving $0$ rewards unchanged.  Clipping the rewards in this manner limits the scale of the error derivatives and makes it easier to use the same learning rate across multiple games.  At the same time, it could affect the performance of our agent since it cannot differentiate between rewards of different magnitude.

In these experiments, we used the RMSProp algorithm with minibatches of size 32.  The behavior policy during training was $\epsilon$-greedy with $\epsilon$ annealed linearly from $1$ to $0.1$ over the first million frames, and fixed at $0.1$ thereafter.  We trained for a total of $10$ million frames and used a replay memory of one million most recent frames.

Following previous approaches to playing Atari games, we also use a simple frame-skipping technique~\cite{bellemare-ale}. More precisely, the agent sees and selects actions on every $k^{th}$ frame instead of every frame, and its last action is repeated on skipped frames.  Since running the emulator forward for one step requires much less computation than having the agent select an action, this technique allows the agent to play roughly $k$ times more games without significantly increasing the runtime.  We use $k=4$ for all games except Space Invaders where we noticed that using $k=4$ makes the lasers invisible because of the period at which they blink.  We used $k=3$ to make the lasers visible and this change was the only difference in hyperparameter values between any of the games.

\subsection{Training and Stability}

%Since there are no convergence guarantees for using Q-learning to learn nonlinear control policies one may wonder whether our proposed approach will work at all. 
In supervised learning, one can easily track the performance of a model during training by evaluating it on the training and validation sets. In reinforcement learning, however, accurately evaluating the progress of an agent during training can be challenging. 
Since our evaluation metric, as suggested by~\cite{bellemare-ale}, is the total reward the agent collects in an episode or game averaged over a number of games, we periodically compute it during training. The average total reward metric tends to be very noisy because small changes to the weights of a policy can lead to large changes in the distribution of states the policy visits .  The leftmost two plots in figure~\ref{fig-V} show how the average total reward evolves during training on the games Seaquest and Breakout.  Both averaged reward plots are indeed quite noisy, giving one the impression that the learning algorithm is not making steady progress.  Another, more stable, metric is the policy's estimated action-value function $Q$, which provides an estimate of how much discounted reward the agent can obtain by following its policy from any given state.  We collect a fixed set of states by running a random policy before training starts and track the average of the maximum\footnote{The maximum for each state is taken over the possible actions.} predicted $Q$ for these states.  The two rightmost plots in figure~\ref{fig-V} show that average predicted $Q$ increases much more smoothly than the average total reward obtained by the agent and plotting the same metrics on the other five games produces similarly smooth curves.
In addition to seeing relatively smooth improvement to predicted $Q$ during training we did not experience any divergence issues in any of our experiments.  This suggests that, despite lacking any theoretical convergence guarantees, our method is able to train large neural networks using a reinforcement learning signal and stochastic gradient descent in a stable manner.  
% --  the predicted values of a held out set of states initially improve quickly before flattening out while the average reward shows a general upward trend but is much noisier. 
%We did not experience any divergence issues in any of our experiments.

%it is quite hard to create such sanitized testing metrics. For that reason, during training we follow the state-value function ($V$) learned by the network and the TD error that Q-learning update minimizes as given in equation~\ref{eq:q-learning}. Both the value $V$ and error is reported over a fixed set of transitions in order to ensure stability in the computation. This can be thought as being similar to a validation set. As shown in figure~\ref{fig-V-TD}, the value predicted by the network increases as the training progresses 
%and only when the $V$ stabilizes, the error starts going down. This kind of behavior shows that as the learning progresses, the agent can play the game better which results in higher expected total reward $V$.

\begin{figure}
\includegraphics[width=0.245\linewidth]{figures/breakout_average_reward}
\includegraphics[width=0.245\linewidth]{figures/seaquest_average_reward}
\includegraphics[width=0.245\linewidth]{figures/breakout_average_v}
\includegraphics[width=0.245\linewidth]{figures/seaquest_average_v}
\caption{\label{fig-V} The two plots on the left show average reward per episode on Breakout and Seaquest respectively during training.  The statistics were computed by running an $\epsilon$-greedy policy with $\epsilon=0.05$ for 10000 steps. The two plots on the right show the average maximum predicted action-value of a held out set of states on Breakout and Seaquest respectively.  One epoch corresponds to 50000 minibatch weight updates or roughly 30 minutes of training time.}
\end{figure}

\subsection{Visualizing the Value Function}
Figure~\ref{fig-reward} shows a visualization of the learned value function on the game Seaquest.  The figure shows that the predicted value jumps after an enemy appears on the left of the screen (point A). The agent then fires a torpedo at the enemy and the predicted value peaks as the torpedo is about to hit the enemy (point B). Finally, the value falls to roughly its original value after the enemy disappears (point C).  Figure~\ref{fig-reward} demonstrates that our method is able to learn how the value function evolves for a reasonably complex sequence of events. 
%Next, we show in figure~\ref{fig-reward} a specific case to demonstrate the behavior of learned value function $V$ on the game of Seaquest. In this game, the player can get rewards by shooting down enemy submarines. From the figure, one can observe that whenever a new submarine appears, the value increases reflecting the fact that the agent has the opportunity to collect more points. As soon as the enemy submarine is shot, the value goes back down since the potential reward is collected.

\begin{figure}
\includegraphics[width=0.24\linewidth]{figures/seaquest_v}
\includegraphics[width=0.24\linewidth]{figures/seaquest_frame_a.png}
\includegraphics[width=0.24\linewidth]{figures/seaquest_frame_b.png}
\includegraphics[width=0.24\linewidth]{figures/seaquest_frame_c.png}
\caption{\label{fig-reward} The leftmost plot shows the predicted value function for a 30 frame segment of the game Seaquest.  The three screenshots correspond to the frames labeled by A, B, and C respectively.}
\end{figure}

\subsection{Main Evaluation}

We compare our results with the best performing methods from the RL literature~\cite{bellemare-ale,bellemare-contingency}.  The method labeled \textbf{Sarsa} used the Sarsa algorithm to learn linear policies on several different feature sets hand-engineered for the Atari task and we report the score for the best performing feature set~\cite{bellemare-ale}.  \textbf{Contingency} used the same basic approach as \textbf{Sarsa} but augmented the feature sets with a learned representation of the parts of the screen that are under the agent's control \cite{bellemare-contingency}.  Note that both of these methods incorporate significant prior knowledge about the visual problem by using background subtraction and treating each of the 128 colors as a separate channel.  Since many of the Atari games use one distinct color for each type of object, treating each color as a separate channel can be similar to producing a separate binary map encoding the presence of each object type.  In contrast, our agents only receive the raw RGB screenshots as input and must \emph{learn} to detect objects on their own.   

In addition to the learned agents, we also report scores for an expert human game player and a policy that selects actions uniformly at random. The human performance is the median reward achieved after around two hours of playing each game. Note that our reported human scores are much higher than the ones in Bellemare et al.~\cite{bellemare-ale}.  For the learned methods, we follow the evaluation strategy used in Bellemare et al. \cite{bellemare-ale,bellemare:qtf} and report the average score obtained by running an $\epsilon$-greedy policy with $\epsilon=0.05$ for a fixed number of steps. The first five rows of table~\ref{table-results} show the per-game average scores on all games.  Our approach (labeled DQN) outperforms the other learning methods by a substantial margin on all seven games despite incorporating almost no prior knowledge about the inputs.   

\begin{table}
\begin{center}
\begin{tabular}{ |l | c | c | c | c | c | c | c |}
	\hline
	& {\small \textbf{B. Rider}} & {\small \textbf{Breakout}} & {\small \textbf{Enduro}} & {\small \textbf{Pong}} & {\small \textbf{Q*bert}} & {\small \textbf{Seaquest}} & {\small \textbf{S. Invaders}}\\
	\hline
	{\small \textbf{Random}}                                    & $354$ & $1.2$ & $0$ & $-20.4$ & $157$ & $110$ & $179$ \\
	\hline
	{\small \textbf{Sarsa}~\cite{bellemare-ale}}                & $996$ & $5.2$ & $129$ & $-19$ & $614$ & $665$ & $271$ \\
	\hline
	{\small \textbf{Contingency}~\cite{bellemare-contingency} } & $1743$ & $6$ & $159$ & $-17$ & $960$ & $723$ & $268$ \\
	\hline
	{\small \textbf{DQN} }                                      & $\textbf{4092}$ & $\textbf{168}$ & $\textbf{470}$ & $\textbf{20}$ & $\textbf{1952}$ & $\textbf{1705}$ & $\textbf{581}$ \\
	\hline
	{\small \textbf{Human} }                                    & $7456$ & $31$ & $368$ & $-3$ & $18900$ & $28010$ & $3690$ \\
	\hline
	\hline
	{\small \textbf{HNeat Best} \cite{hausknecht-neuro}}        & $3616$ & $52$ & $106$ & $19$ & $1800$ & $920$ & $\textbf{1720}$ \\
	\hline
	{\small \textbf{HNeat Pixel} \cite{hausknecht-neuro}}       & $1332$ & $4$ & $91$ & $-16$ & $1325$ & $800$ & $1145$ \\
	\hline
	{\small \textbf{DQN Best}}                                  & $\textbf{5184}$ & $\textbf{225}$ & $\textbf{661}$ & $\textbf{21}$ & $\textbf{4500}$ & $\textbf{1740}$ & $1075$ \\
	\hline
\end{tabular}
\caption{\label{table-results} The upper table compares average total reward for various learning methods by running an $\epsilon$-greedy policy with $\epsilon=0.05$ for a fixed number of steps. The lower table reports results of the single best performing episode for HNeat and DQN.  HNeat produces deterministic policies that always get the same score while DQN used an $\epsilon$-greedy policy with $\epsilon=0.05$.}
\end{center}
\end{table}

We also include a comparison to the evolutionary policy search approach from~\cite{hausknecht-neuro} in the last three rows of table~\ref{table-results}.   
%Since only the maximum score obtained on each game was reported for this method, we compare it to the maximum score for our method.  
We report two sets of results for this method. The \textbf{HNeat Best} score reflects the results obtained by using a hand-engineered object detector algorithm that outputs the locations and types of objects on the Atari screen. The \textbf{HNeat Pixel} score is obtained by using the special 8 color channel representation of the Atari emulator that represents an object label map at each channel. 
This method relies heavily on finding a deterministic sequence of states that represents a successful exploit. It is unlikely that strategies learnt in this way will generalize to random perturbations; therefore the algorithm was only evaluated on the highest scoring single episode. In contrast, our algorithm is evaluated on $\epsilon$-greedy control sequences, and must therefore generalize across a wide variety of possible situations.  Nevertheless, we show that on all the games, except Space Invaders, not only our max evaluation results (row $8$), but also our average results (row $4$) achieve better performance.
%More importantly, this method relies heavily on finding a unique sequence of states that represent a successful exploit. Therefore, using a stochastic evaluation procedure which avoids overfitting is not suitable with this approach and the results are reported by searching for the single best episode that the algorithm plays. Despite these aspects, we show that on all the games, except Space Invaders, not only our max evaluation results (row $8$), but also our average results (row $4$) achieve better performance.  

Finally, we show that our method achieves better performance than an expert human player on Breakout, Enduro and Pong and it achieves close to human performance on Beam Rider. The games Q*bert, Seaquest, Space Invaders, on which we are far from human performance, are more challenging because they require the network to find a strategy that extends over long time scales.  
